\section{Fresh finite-contact entry and the exact two-family source target}
\label{sec:rederived-contact-entry}

The input here is the actual recovered parent and its factored determinant,
not an asserted completion of an atlas.  We rederive the finite-contact
nondegeneracy needed by the source-entry argument and state exactly which
C1/C2 certificates would finish the global source reduction.  A strict
Hessian in a single source field is not identified with a strict real
residue Gram in a tensor product.

\subsection{Exact input and support incidence}

Use the full parent (CAP1), with $d>0$ and all eight event coefficients
nonnegative.  Its state matrix is $P(u,v)$.  The physical dictionary is
$F=D_uD_v\|\psi\|^2 f(q)$ at a nonzero state and its homogeneous
extension at projective angles.  This dictionary and the determinant
identity below occur as actual proof bodies in the recovered
\emph{physical\_core.tex}; the physical normalization was also freshly
rederived in (RFA10)--(RFA12).

Write
\begin{align*}
t_1&=d+2(\lambda+n+D),&t_2&=d+2(\alpha+m+A),\\
t_3&=d+2(\lambda+C+D),&t_4&=d+2(\alpha+\beta+m),\\
t_5&=d+2(\lambda+\beta+n),&t_6&=d+2(\alpha+C+A),\\
t_7&=d+2(\beta+n+m),&t_8&=d+2(C+A+D).
\end{align*}
All $t_j$ are strictly positive.  The exact whole-parent determinant is
\[
\det P=16D_u^2D_v^2u^2v^2L(u^2,v^2),\qquad
L(U,V)=h+\sum_{e\in E}q_eU^{e_1}V^{e_2},
\quad E=\{-1,0,1\}^2\setminus\{0\},
\tag{SE1}
\]
where $h=q_{11}$ and the eight nonnegative coefficients are
\begin{align*}
q_{00}&=4\lambda\alpha t_7t_8,&
q_{01}&=4(\lambda Dt_7t_6+\alpha mt_5t_8),\\
q_{02}&=4mDt_5t_6,&
q_{10}&=4(\alpha At_7t_3+\lambda nt_8t_4),\\
q_{12}&=4(\beta mt_1t_6+CDt_5t_2),&
q_{20}&=4nAt_3t_4,\\
q_{21}&=4(\beta nt_3t_2+CAt_1t_4),&
q_{22}&=4\beta Ct_1t_2.
\end{align*}
Here $q_{ij}$ has exponent $e=(i-1,j-1)$.  In particular an exponent
is active exactly when at least one of its displayed event products is
positive.  The checker for this section uses these exact incidence
conditions; it does not replace physicality by a support enumeration.

Call the support \emph{coercive} when $0$ is in the interior of the
convex hull of the active exponents.  Over the real field this is
equivalent to coercivity of $L(e^x,e^y)$: a positive exponential
sum tends to infinity in every direction exactly under that condition.
The condition on the finite support itself makes sense unchanged over
any real closed coefficient field.

At $Z=0$ the signed part of the full parent is
\[
2d\{(uW+uvX-Y)^2+D_uv^2Y^2+D_vX^2\}.
\tag{SE2}
\]
It is positive definite for $uv\ne0$, so the corresponding
$3\times3$ block $H$ is positive definite, and
$P\succeq0\Longleftrightarrow\det P\geq0$.  If a nonzero state is
in its kernel, that state has $Z\ne0$ and the kernel has dimension one.

\subsection{No physical finite zero has degenerate Newton support}

\begin{theorem}[Coercivity and strict curvature at every actual finite zero]
Let a nonnegative canonical source with $d>0$ be physical and have a
zero on a nonzero state at finite $u_c,v_c\ne0$.  Then its Newton
support is coercive.  After the harmless real local basis gauges making
$u_c,v_c>0$ and the normalization $Z_c=1$, the full Hessian of $F$ in
the five variables $(W,X,Y,u,v)$ is positive definite at this contact.
The assertion holds over every real closed extension of $\mathbb R$.
\end{theorem}

\begin{proof}
Physicality makes the contact a zero minimum, so its first derivatives
in the five variables vanish.  By (SE1)--(SE2), $L=0$ and its two
logarithmic derivatives vanish.  Put
$a_e=q_eU_c^{e_1}V_c^{e_2}>0$ for active exponents.  Then
\[
\sum_ea_e e=0.
\tag{SE3}
\]
If the support were not coercive, a nonzero real linear functional
would be nonnegative on it.  Applying that functional to (SE3) shows
that it vanishes on every active exponent.  Thus all active exponents
would lie on one line through the origin, or the support would be empty.
There are only four possible such lines in $E$.

An active horizontal opposite pair would require both
$q_{01},q_{21}>0$.  Their four possible product pairings force,
respectively, $\lambda n$, $CD$, $\beta m$, or $\alpha A$ to be
positive, activating $q_{10}$ or $q_{12}$ off that line.  A vertical
opposite pair $q_{10},q_{12}>0$ is ruled out in exactly the same
explicit way: its pairings force $\alpha m$, $CA$, $\beta n$, or
$\lambda D$ positive, activating $q_{01}$ or $q_{21}$.  A
one-sided line cannot satisfy (SE3).

For the diagonal pair $q_{00},q_{22}>0$, positivity forces
$\lambda,\alpha,\beta,C>0$, and all off-line coefficients vanish
only if $n=m=A=D=0$.  For the other diagonal pair
$q_{02},q_{20}>0$, the sole possibility is
$n,m,A,D>0$, $\lambda=C=\alpha=\beta=0$.
These are the two alternating four-event matchings of the cycle.
Neither supports a physical source with $d>0$.  For the latter, set
$u=v=1$ and take the entire, untruncated state
$\psi=(1,-\varepsilon,\varepsilon,1)$.  Direct substitution gives
\[
F(1,1;1,-\varepsilon,\varepsilon,1)
 =16\varepsilon\{(d+n+m+A+D)\varepsilon-d\}.
\tag{SE4}
\]
Choosing $\varepsilon=d/[2(d+n+m+A+D)]>0$ is a strict negative
witness.  The other matching is its image under the physical one-step
event-cycle rotation (CAP13).  The formula also allows literal zeros
among the matching weights.

Finally suppose every outer coefficient vanishes.  In the event cycle
$(\lambda,D,C,A,\alpha,m,\beta,n)$, the forbidden simultaneously
positive pairs are precisely the adjacent pairs and opposite pairs:
these are the twelve products occurring in (SE1).  Four pairwise
nonadjacent vertices of an eight-cycle would have to alternate and
would contain an opposite pair.  Hence at most three events are active.
Every maximal allowed set is a triple with cyclic gaps $(2,3,3)$;
smaller allowed sets extend to one.  These eight triples are all rotations
of $\{\lambda,m,A\}$.  At $u=v=1$, the state $(1,-1,0,2)$
annihilates those three event rows, while the signed core equals $-4$.
Thus the full parent is $-4d<0$.  Rotation rules out every such support.
This excludes the empty and rank-one possibilities, proving coercivity.

The angular log-Hessian at the contact is now
\[
\mathcal H_L=\sum_ea_e ee^T\succ0,
\tag{SE5}
\]
because the active exponent vectors span $\mathbb R^2$ and each
$a_e$ is positive.  Write
$P=\left(\begin{smallmatrix}H&b\\b^T&c\end{smallmatrix}\right)$.
Completing the three state coordinates within the single source field
gives
\[
F=(x+H^{-1}b)^TH(x+H^{-1}b)+s(u,v),\qquad
s=c-b^TH^{-1}b=\det P/\det H.
\tag{SE6}
\]
At the contact $s=0$ and $\nabla s=0$.  Equations (SE1), (SE5)
show that its angle Hessian is the positive scalar
$16D_u^2D_v^2u^2v^2/\det H$ times the congruence of
$\mathcal H_L$ by $\operatorname{diag}(2/u,2/v)$.
It is positive definite.  The full five-variable Hessian has positive
state block $2H$ and this positive angular Schur complement, so is
positive definite too.  This is an exact single-field Schur argument.
Logarithms are merely notation for the operators
$U\partial_U,V\partial_V$; (SE3), (SE5), and their Hessian
congruence are algebraic ordered-field formulas.  Alternatively the
entire fixed finite assertion transfers from the real field.
\end{proof}

There is a useful converse contact implication, independent of
coercivity.  If nonnegative coefficients with $d>0$ have a nonzero
kernel state at $uv\ne0$ and both angle derivatives of its parent
vanish, (SE2) gives $L=0$, $U\partial_UL=V\partial_VL=0$.
Over $\mathbb R$, convexity of $L(e^x,e^y)$ makes this zero its
global minimum.  The determinant criterion and the direct commuting
endpoint argument prove physicality.  This is a finite first-order
implication, so it transfers to every real closed field.  Thus
\[
\text{nonnegative coefficients + finite kernel + both angle jets}
 \quad\Longrightarrow\quad\text{physicality}.
\tag{SE7}
\]
No coercivity premise is needed for this implication.  Combining it
with the theorem then excludes every proposed degenerate-support
contact automatically.

\subsection{The negative-core state orientation and deletion test}

For reference, the raw signed-core matrix $P_H$ has
\[
\det P_H=-16D_u^2D_v^2u^2v^2\ne0.
\tag{SE8}
\]
At any physical finite zero, at least one retained event has positive
value.  Otherwise all event matrices annihilate the state and the
kernel equation would give $P_H\psi=0$, contradicting (SE8).
Consequently the core expectation there is strictly negative.

This fixes an exact state sign without native contact coordinates.
For $u,v>0$ put $h_0=1+v^2+u^2v^2$.  The $X,Y$ core block is
$2\left(\begin{smallmatrix}h_0&-uv\\-uv&h_0\end{smallmatrix}\right)
\succ0$.  Its Schur complement on $(W,Z)$ is
\[
\frac{2D_uD_v}{h_0^2-u^2v^2}
\begin{pmatrix}u^2v^2&-uvh_0\\-uvh_0&u^2v^2\end{pmatrix}.
\tag{SE9}
\]
If $WZ\leq0$, this quadratic form and the completed $X,Y$ square
are nonnegative.  Hence a negative core expectation forces $WZ>0$.
At a physical contact the gauge $Z=1$ therefore has $W>0$.
For an independent testing state with $WZ\leq0$, the core value is
instead a nonnegative pure-testing scalar; its separated product with
$d$ is a square, and every event is a square.  This gives an exact
easy-state tensor-SOS exit, with no claim about merely small positive
$WZ$.

Suppose a positive retained event has zero row at a finite contact.
Deleting its coefficient preserves the kernel and both angle jets:
the square and all its first derivatives vanish there.  The remaining
coefficients are nonnegative, so (SE7) makes the deleted source physical.
The theorem therefore forbids a deletion that makes the Newton support
noncoercive.  This supplies the exact no-wall implication used in the
recovered global C1 sign-entry proof.  It does not authorize deleting
an event whose row is nonzero: (CAP15) then gives a negative value.
\tag{SE10}

\subsection{The canonical event basis and two sufficient contact families}

Here we specify the global source target, including its quantifiers.
Let $Q\subset\mathbb R^8$ be the compact finite-dimensional
tensor-product quantum set from (RFA1)--(RFA6), and let $f$ be an affine
functional nonnegative on it.  If $f$ is nonnegative on the whole
nonsignaling polytope, linear separation of its finite probability
inequality description writes it as a nonnegative sum of probability
events (a nonnegative constant can be absorbed using normalization).
These events have immediate separated certificates.  Otherwise $f$ is
negative at a PR vertex: the nonsignaling vertices are the sixteen
deterministic assignments and eight PR vertices, and $f$ is nonnegative
at every deterministic assignment.

Choose that PR orientation.  In the convention
$H=1-(E_{00}-E_{01}+E_{10}+E_{11})/2$, the eight losing events
vanish at that PR point.  Each has a unique deterministic $S=2$
neighbor where its probability is one and the seven other losing
probabilities vanish.  The nine evaluations are therefore an affine
basis and give the exact coefficient representation
\[
f=dH+4\sum_i a_i p_i,\qquad
d=-f(\mathrm{PR})>0,\qquad a_i=f(v_i)/4\geq0.
\tag{SE11}
\]
The checker verifies these eight neighbor evaluations and nonsingularity
of their augmented nine-point evaluation matrix from the actual event
definitions.  The finite nonsignaling vertex argument and the full
physical normalization are retained inputs, with exact source hashes
in the accompanying dependency receipt.

Two omissions are compatible precisely when the events never specify
opposite outcomes at a shared local setting.  The factored coefficients
in (SE1) give
\[
0\in\operatorname{int}\operatorname{conv}(\operatorname{supp}L)
\quad\Longleftrightarrow\quad
\text{all omitted losing events are pairwise compatible}.
\tag{SE12}
\]
For clarity, a separating direction for a subset of the eight square
boundary points can be chosen from
$(\pm1,0),(0,\pm1),(\pm1,\pm1)$.  For an axis representative,
vanishing on the forbidden half-plane requires
$\lambda\alpha=\lambda D=\alpha m=mD=0$, whose minimal
omission sets are $\{\lambda,m\}$ and $\{\alpha,D\}$.
For a diagonal representative the minimal sets are
$\{\lambda,\alpha\}$, $\{\lambda,m,A\}$,
$\{n,D,\alpha\}$; the latter two already contain an axis
exclusive pair.  The physical cycle rotation gives the other directions.
These are exactly the twelve exclusive pairs: cycle distance three or
four.  This proves (SE12), and the checker independently verifies it on
all finite support choices.  Compatible pairs have cycle distance one
or two; a compatible triple is three consecutive vertices.  A larger
compatible omission set is impossible.

Normalize the nonnegative canonical source tuple by
$d+\sum a_i=1$.  Its physical subset is compact and convex over
$\mathbb R$.  Consider an extreme point with $d>0$ and coercive
support.  Its Laurent polynomial attains a minimum $\mu$ at finite
positive $U,V$.  If $\mu>0$, perturb the active coefficients by
small opposite nonradial amounts, preserving their sum and exact zero
coordinates.  Every positive $q_e$ changes by relative error at most
$\epsilon$ and $h$ by absolute error at most $\epsilon$.  Uniformly,
\[
L'\geq(1-\epsilon)L+\epsilon h-\epsilon
 \geq(1-\epsilon)\mu+\epsilon h-\epsilon>0
\tag{SE13}
\]
for sufficiently small real $\epsilon$.  The perturbed sources remain
physical, contradicting extremality.  There are at least six active
coordinates because (SE12) permits at most three event omissions, so
a nonradial perturbation exists.  Hence $\mu=0$, giving an actual
finite noncommuting contact.

At that contact the four kernel equations and two angle equations
\[
P\psi=0,\qquad \psi^TP_u\psi=0,\qquad\psi^TP_v\psi=0
\tag{SE14}
\]
are six homogeneous linear equations in the nine source coefficients.
If seven or eight events were positive, there would be at least eight
active source coordinates including $d$.  Their nullspace has dimension
at least two and contains the source itself.  It therefore has a
nonzero direction of coefficient sum zero.  Small opposite changes in
that direction preserve all active signs and the same equations
(SE14).  By (SE7) they remain physical, again contradicting extremality.
No generic-rank assumption is involved.  Thus a coercive extreme source
has exactly five or six positive events.

Choose an adjacent compatible omitted pair first when three events are
omitted.  By the physical cycle symmetries, every such source enters
one of the following exact domains, including the stated zero faces:
\[
\begin{array}{ll}
\mathrm{C1}:&d,m,C,A,\alpha>0,\quad
 \lambda=n=0,\quad D,\beta\geq0,\quad D+\beta>0,\\
\mathrm{C2}:&d,m,C,A,D,\alpha>0,\quad
 \lambda=\beta=0,\quad n\geq0.
\end{array}
\tag{SE15}
\]
Both domains retain physicality and the existence of an actual finite
positive source contact.  In C1 a third omitted event can only be
$D$ or $\beta$.  In C2 it can only be $n$; this boundary is already
in a C1 orbit, but allowing it in both is harmless.  Noncoercive source
faces contain an exclusive omitted pair and are relabelings of
$\lambda=m=0$ or $\lambda=\alpha=0$; their recovered exact matrix
lifts remain separate, available proof inputs.

The normalized source base has affine dimension at most eight, so the
finite-dimensional extreme-point theorem and Caratheodory's elimination
argument express each point as a convex combination of at most nine
extreme points.  The midpoint definition of extremality, physicality,
this fixed nine-point combination, and finite-contact attainment all
have first-order polynomial formulas.  Their truth over $\mathbb R$
therefore transfers to every real closed coefficient field.  Compactness,
logarithms, and real perturbation sizes were used only to prove the real
assertion before transfer.  Nonnegative summands automatically preserve
literal zeros of the original source.

\begin{corollary}[Precise sufficient tensor hypotheses]
Suppose that, for each of C1 and C2 in (SE15), every actual source
tuple over every real closed $K_c\supset\mathbb R$ has an
unmultiplied full-parent SOS at every independent projective test and
state over every real closed $K_t\supset\mathbb R$.  There is no
algebraic-disjointness assumption: the fields are labelled factors and
may be copies of the same field.  Together with the two recovered
closed-face matrix lifts, these two hypotheses imply the universal
separated slack-SOS statement for the entire physical set $Q$.
\end{corollary}

\begin{proof}
Apply (SE11)--(SE15) in the source field, restore the finite real Bell
relabelings, and multiply roots by the separate square roots of the
nonnegative source combination weights.  The event atoms and the
closed-face atoms have separated certificates, the latter by attained
finite-dimensional matrix duality as in (RFA).  Every quantum testing
point is a convex combination of at most nine real-qubit atoms in its
own field; multiply roots by the separate testing-weight square roots
and sum.  The pure homogeneous angle, state and source normalization
factors are restored within their respective fields.  These operations
are exactly finite conic and convex closure of sums of squares.
\end{proof}

\paragraph{What is and is not discharged.}
This section proves source-contact coercivity, the strict single-field
contact Hessian, contact sufficiency, the no-wall deletion implication,
and the exact canonical source partition.  It isolates sufficient C1/C2
tensor hypotheses with every coefficient sign, literal boundary and
source contact retained.  Their positive-total and weak-total subdomains
are the exhaustive scale split (CAP17).  It does not discharge those
two all-tests hypotheses.  In particular, (SE6) is not a mixed tensor
inverse, and its positive Hessian need not have a positive definite
real residue after arbitrary independent source/testing degenerations.
No generic matrix descent or historical all-tests assertion is used.
